\ifnotes
\clearpage
\section*{Appendix: ACS requirements for this document}
\addcontentsline{toc}{section}{Appendix}
\small

Internal checklist, compiled from the ACS General Grant Policies and the General
Grant Application Instructions, both effective September 2026. Not part of the
submitted proposal: set \texttt{\textbackslash notesfalse} in the preamble to
drop this appendix before export. Page numbering restarts nowhere, so removing
it leaves the body pagination unchanged.

\subsection*{A. The research plan itself (Section 6.1)}

Nine pages maximum, excluding references. The plan must be realistic for the
period of support requested, and a timeline is encouraged. Six required
elements, in this order:

\begin{enumerate}
\item \textbf{Specific Aims.} List the objectives and goal of the research
  proposed, and describe the specific aims briefly, in order of priority.
\item \textbf{Background and Significance.} Concisely summarise and critically
  evaluate relevant work done by others that provides a rationale or
  justification for the project. State specifically how successful completion
  of the work will advance scientific knowledge relevant to cancer initiation,
  prevention, detection, treatment or survivorship.
\item \textbf{Preliminary Studies.} Optional. If included, give results of
  research accomplished by the applicant or others that is relevant to the
  proposal, explain its significance to the scientific rationale, and
  \emph{carefully attribute the source of any preliminary data}.
\item \textbf{Research Design and Methods.} Methods and procedures in
  sufficient detail to permit evaluation by other scientists. Must include: a
  statistical analysis plan; a power analysis to support design and sample size
  where applicable; expected outcomes and anticipated benchmarks for success;
  potential difficulties and limitations of the methods; and alternative
  methods and approaches.
\item \textbf{Impact.} How the anticipated outcomes will affect concepts,
  methods, technologies, treatments, services or interventions, and move the
  field forward.
\item \textbf{References.} Numbered in order of appearance in the text. Each
  must include title, names of \emph{all} authors, book or journal, volume
  number, page numbers and year. Not counted against the page limit.
\end{enumerate}

\subsection*{B. Formatting, applying to every uploaded template}

\begin{itemize}
\item Principal Investigator name in the header of each template. Do not change
  the footers on ACS-supplied templates.
\item Single or double spaced. If single spaced, leave a blank space between
  paragraphs.
\item Type size: 12-point Times New Roman or 11-point Arial minimum for text;
  10-point Times New Roman or 9-point Arial may be used for figures, legends
  and tables.
\item Margins greater than 0.5 inches on all sides, unless a supplied template
  sets different margins.
\item Page numbers in the upper right corner, according to the proposal
  sections listed in the Table of Contents.
\item Do not number: Signature Page, Contact Page, General Audience Summary,
  Structured Technical Abstract, Statement of Cancer Relevancy and Impact,
  Justification of Alignment with Research Priorities, Budget and
  Justification, and the Appendix.
\item All attachments, including the Appendix, upload as PDF under
  ``Application Sections'' in ProposalCentral, each with the correct attachment
  type selected.
\end{itemize}

\subsection*{C. What reviewers score this document against}

\begin{itemize}
\item \textbf{Research plan summary, significance and cancer relevance.}
  Whether the project addresses an important problem or critical barrier and is
  important to cancer research; how completing the aims would change concepts,
  methods, technologies, treatments, services or interventions; how the work is
  relevant to people at risk for or living with cancer, their families and
  caregivers, or populations facing an unequal burden. Relevance to cancer may
  be indirect, but the connection must be clearly articulated by the applicant.
\item \textbf{Research plan approach.} Whether the hypothesis, objectives and
  aims are appropriate for the project period; whether the strategy,
  methodology and analyses are well reasoned; the quality of any preliminary
  data, judged for feasibility and justification; whether statistical
  considerations, potential problems, alternative strategies and benchmarks for
  success are clearly presented. For early-stage or hypothesis-generating work,
  whether the study is well justified and the risks mitigated. A postdoctoral
  fellow's plan is not expected to show the breadth and depth of a senior
  scientist's, but it must be fundamentally sound.
\item \textbf{Candidate.} Whether the proposed training and research
  \emph{broadens the applicant's experience beyond their graduate work} and
  aligns with their stated career goals.
\end{itemize}

\subsection*{D. The rest of the package, for reference}

Applicant's sections: Table of Contents (2 pages max, with appendices itemised
at the bottom); Pending Fellowship Applications (2.1); Biographical Sketch on
the ACS template, 5 pages max (3.1); Reply to Previous Review (4.1), marked
``Not Applicable'' for a new submission; Statement of Experience, Training
Potential and Career Goals, 3 pages max (5.1); this research plan (6.1); and
Letters of Recommendation from three people other than the mentor (7.1), which
must be uploaded before the application can be submitted.

Mentor's sections, prepared by the primary mentor: Proposed Training and
Mentoring, 3 pages max (8.1), which must state whether the research plan was
prepared independently by the applicant or in collaboration with the mentor;
Mentor Commitment Letter (9.1); Research and Training Environment, 3 pages max
(10.1); NIH-format Biosketch via SciENcv (11.1); Research Support of Mentor
(12.1), listing granting agency, project title, direct costs with per-year or
total clearly indicated, and term of award; and Compliance Statements (13.1).

Typed into ProposalCentral rather than uploaded: General Audience Summary
(3,100 characters, lay language, scored by Community Research Partners);
Structured Technical Abstract (3,100 characters, organised as Background,
Objective/Hypothesis, Specific Aims, Study Design); and Statement of Cancer
Relevance and Impact (1,500 characters, no technical jargon).

\subsection*{E. Terms}

Three-year fellowship: \$217,500 total. Stipend \$66,000, \$68,000 and \$70,000
across years one to three, salary only and not intended to cover fringe
benefits, though the institution may supplement it provided 100\% research
effort on the project is maintained. Allowance \$4,000 per year, plus \$1,500
in the final year prioritised for travel to the ACS Jiler Professors and
Fellows Conference. No indirect costs. The fellow must remain in a mentored,
non-faculty role for the duration of the award, and may not hold another
postdoctoral fellowship concurrently. ACS does not fund projects that share
specific aims with other funded or pending work.

\subsection*{F. Use of AI tools}

The applicant is responsible for all content in the application, including
anything generated with an AI tool or large language model. ACS asks that such
tools be credited and sources cited where possible, and that institutional
guidelines be followed.
\fi